\documentclass[11pt,a4paper]{article}
\usepackage{amsmath, amssymb, geometry}
\usepackage{graphicx}
\geometry{margin=2.5cm}
\setlength{\parskip}{0.9em}
\setlength{\parindent}{0pt}

\begin{document}

\section*{Negative Log-Likelihood and its Lower Bound}

Let $C \in \{L, C, R\}$ denote the subject's choice and $S$ the stimulus or condition.
A stochastic model with parameters $\theta$ defines conditional probabilities
$q_\theta(c \mid s)$.
Given data $\{(c_i, s_i)\}_{i=1}^N$, the mean negative log-likelihood (NLL)
per trial is

\[
\overline{\mathrm{NLL}}(\theta)
= -\frac{1}{N}\sum_{i=1}^N \log q_\theta(c_i \mid s_i).
\]

Under the i.i.d.\ assumption, the empirical mean converges to its expected value:

\[
\overline{\mathrm{NLL}}(\theta)
\xrightarrow[]{N\to\infty}
-\mathbb{E}_{p(C,S)}[\log q_\theta(C \mid S)]
= H(p(C \mid S), q_\theta(C \mid S)),
\]

which is the \emph{cross-entropy} between the subject’s true conditional
distribution $p(C \mid S)$ and the model’s predicted probabilities
$q_\theta(C \mid S)$.

According to Cover and Thomas (2006, Sec.~2.6), the cross-entropy decomposes as

\[
H(p(C \mid S), q_\theta(C \mid S))
= H(C \mid S)
+ \mathbb{E}_{s\sim p}\!\left[
D_{\mathrm{KL}}(p(\cdot \mid s)\, \|\, q_\theta(\cdot \mid s))
\right],
\]

where $D_{\mathrm{KL}}$ denotes the Kullback–Leibler divergence
Since $D_{\mathrm{KL}} \ge 0$, it follows that

\[
\boxed{\overline{\mathrm{NLL}}(\theta) \ge H(C \mid S)}
\qquad
\text{with equality iff } q_\theta(c \mid s) = p(c \mid s) \;\forall c,s.
\]

The lower bound $H(C \mid S)$, the \emph{conditional entropy}, is defined as

\[
H(C \mid S)
= -\sum_s p(s) \sum_{c\in\{L,C,R\}} p(c \mid s) \log p(c \mid s),
\]
measured in nats.

To estimate $H(C \mid S)$ empirically, let $n_s$ be the number of trials with
stimulus $s$, $n_{c \mid s}$ the number of times the subject chose $c$
given $s$, and $N = \sum_s n_s$.
Empirical frequencies are
\[
\hat{p}(s) = \frac{n_s}{N}, \qquad
\hat{p}(c \mid s) = \frac{n_{c \mid s}}{n_s} \quad (n_s > 0).
\]
The plug-in estimator of the conditional entropy is then
\[
\widehat{H}(C \mid S)
= -\sum_s \hat{p}(s)
  \sum_c \hat{p}(c \mid s)\, \log \hat{p}(c \mid s),
\]
where the convention $0\log0 \equiv 0$ is used.

For $K=3$ equiprobable alternatives (chance level), the reference NLL is
\[
\mathrm{NLL}_{\text{chance}} = \log K = \log 3 \approx 1.099
\ \text{nats per trial.}
\]
The subject’s intrinsic noise limits the explainable information, determined by
$H(C \mid S)$.

A normalized fit metric can be defined as
\[
\mathrm{Fit}_{\text{ceiling}}
= 1 -
\frac{\overline{\mathrm{NLL}}(\theta) - \widehat{H}(C \mid S)}
{\log K - \widehat{H}(C \mid S)} \in [0,1],
\]
where 0 corresponds to chance level and 1 corresponds to a perfect fit
to the subject’s behavior.

If $q_\theta(c \mid s)$ is estimated via Monte Carlo simulation with $M$
trajectories (and smoothing $\alpha$), the empirical probability used in the loss is
$\widehat{q}(c \mid s) = \frac{m_c+\alpha}{M+K\alpha}$.
This introduces a small positive bias in the NLL;
in the limit $M\!\to\!\infty$ and $\alpha\!\to\!0$,
$\overline{\mathrm{NLL}}(\theta)$ converges to the theoretical bound above.


\subsection*{Monte Carlo Estimation of the Cross-Entropy}

The theoretical cross-entropy between the true behavior $p(c\mid s)$ and the
model prediction $q_\theta(c\mid s)$ is
\[
H(p,q_\theta)
= -\mathbb{E}_{(c,s)\sim p}\,[\log q_\theta(c\mid s)].
\]
When $q_\theta(c\mid s)$ cannot be computed analytically, it can be estimated
via Monte Carlo simulation with $M$ stochastic trajectories per condition:
\[
\widehat{q}_\theta(c\mid s)
= \frac{m_c + \alpha}{M + K\alpha},
\]
where $m_c$ is the number of simulated choices equal to $c$,
$K$ is the number of possible choices, and $\alpha$ is a small smoothing term.
Replacing $q_\theta$ by $\widehat q_\theta$ in the definition yields the
Monte Carlo estimator of the cross-entropy:
\[
\widehat{H}_{MC}(p,q_\theta)
= -\frac{1}{N}\sum_{i=1}^N \log \widehat q_\theta(c_i\mid s_i)
\approx H(p,q_\theta).
\]
As $M\!\to\!\infty$, $\widehat q_\theta(c\mid s)\!\to\! q_\theta(c\mid s)$ and
$\widehat{H}_{MC}(p,q_\theta)\!\to\! H(p,q_\theta)$.
For finite $M$, the estimator is slightly biased upwards due to Jensen's
inequality, with bias $\mathcal{O}(1/M)$
(MacKay, 2003; Bishop, 2006; Cover \& Thomas, 2006).

\subsection*{Effect of Subsampling and Reweighting}

In practice, behavioral datasets are often subsampled or balanced
across conditions to obtain approximately uniform trial counts.
Let $\tilde p(s)$ denote the empirical frequency of each condition $s$
after subsampling.
Then the estimated cross-entropy becomes
\[
\widehat{H}^{(\text{sub})}_{MC}(p,q_\theta)
= -\sum_s \tilde p(s)\!
   \sum_c \hat p(c\mid s)\, \log \widehat q_\theta(c\mid s),
\]
where $\hat p(c\mid s)$ are the observed conditional frequencies.
Subsampling therefore modifies the weighting $\tilde p(s)$
but not the conditional terms $\hat p(c\mid s)$.
As long as the same weighting is used when computing both
$\overline{\mathrm{NLL}}(\theta)$ and $\widehat H(C\mid S)$,
the inequality
\[
\overline{\mathrm{NLL}}(\theta)
\;\ge\; H(C\mid S)
\]
remains valid up to sampling noise.
Different subsampling schemes simply change the effective mixture of
conditions $s$ over which the cross-entropy is averaged,
but do not alter the theoretical relationship between
the NLL and the conditional entropy bound.

\vspace{1em}
\textbf{References}

Cover, T.\,M. \& Thomas, J.\,A. (2006).
\emph{Elements of Information Theory} (2nd ed.).
Wiley.\\
Bishop, C.\,M. (2006).
\emph{Pattern Recognition and Machine Learning}.
Springer.\\
MacKay, D.\,J.\,C. (2003).
\emph{Information Theory, Inference, and Learning Algorithms}.
Cambridge University Press.

\end{document}